% ----------------------------------------------------------------------
%  2.2.3. Упругие свойства дислокаций
% ----------------------------------------------------------------------
\subsection{Упругие свойства дислокаций}\label{ssec:disl-elastic}

\begin{orientir}
Дислокация искажает решётку вокруг себя на большие расстояния. Это поле
напряжений определяет почти всё её «поведение»: энергию дислокации,
силы, с которыми на неё действуют внешние нагрузки и другие дефекты,
притяжение примесных атомов. План раздела: поля напряжений винтовой и
краевой дислокаций $\to$ энергия дислокации $\to$ силы на дислокацию
(формула Пича--Кёлера) $\to$ взаимодействие дислокаций между собой и с
точечными дефектами (атмосферы).
\end{orientir}

\subsubsection{Ядро и упругая область}

По определению дислокация создаёт поле смещений $\vb u(\vb r)$, для
которого на любом контуре вокруг линии $\oint\d\vb u=\vb b$. Искажения
вокруг дислокации делят на две области:
\begin{itemize}
\item \term{ядро} ($r<r_0$, $r_0\approx(2\text{--}3)\,b$): смещения
атомов сравнимы с межатомным расстоянием, законы линейной упругости не
работают;
\item \term{упругая область} ($r>r_0$): деформации малы, и решётку можно
считать сплошной упругой средой.
\end{itemize}
Мы будем описывать упругую область, где можно~— отмечая вклад ядра.

Нам понадобятся два тензора (см. п.~1.5). \term{Тензор деформаций}
\[
\varepsilon_{ij}=\frac12\left(\frac{\partial u_i}{\partial x_j}+\frac{\partial u_j}{\partial x_i}\right)
\]
описывает, как смещаются точки тела. \term{Тензор напряжений}
$\sigma_{ij}$~— это $i$-я компонента силы, действующей на единичную
площадку с нормалью вдоль оси $j$ (рис.~\ref{fig:stress-comp}). Для
изотропной среды они связаны законом Гука:
\begin{equation}\label{eq:hooke}
\sigma_{ij}=2G\varepsilon_{ij}\quad(i\neq j),\qquad
\sigma_{ii}=2G\left(\varepsilon_{ii}+\frac{\nu}{1-2\nu}\,\varepsilon_{ll}\right),
\end{equation}
где $\varepsilon_{ll}=\varepsilon_{11}+\varepsilon_{22}+\varepsilon_{33}$~—
дилатация (относительное изменение объёма), $\nu$~— коэффициент Пуассона.

\begin{figure}[ht]
\centering
\begin{subfigure}[b]{0.45\textwidth}\centering
\begin{tikzpicture}[cubeproj, scale=1.6]
  \CubeEdges
  \draw[arr, fmRed] (1,0.5,0.5) -- (1.6,0.5,0.5) node[right, slbl]{$\sigma_{xx}$};
  \draw[arr, fmBlue] (1,0.5,0.5) -- (1,1.15,0.5) node[right, slbl]{$\sigma_{yx}$};
  \draw[arr, fmBlue] (1,0.5,0.5) -- (1,0.5,1.0) node[above, slbl]{$\sigma_{zx}$};
  \draw[arr] (0,0,0) -- (1.9,0,0) node[right, slbl]{$x$};
  \draw[arr] (0,0,0) -- (0,0,1.5) node[above, slbl]{$z$};
  \draw[arr] (0,1,0) -- (0,1.7,0) node[right, slbl]{$y$};
\end{tikzpicture}
\caption{}
\end{subfigure}\hfill
\begin{subfigure}[b]{0.45\textwidth}\centering
\begin{tikzpicture}[scale=1.1]
  \draw[arr] (0,0) -- (2.6,0) node[right, slbl]{$x$};
  \draw[arr] (0,0) -- (0,2.3) node[above, slbl]{$y$};
  \draw[fmBlue!60] (0,0) -- (40:2.3);
  \draw (0.6,0) arc (0:40:0.6); \node[slbl] at (20:0.85) {$\theta$};
  \coordinate (P) at (40:1.7);
  \fill (P) circle (1.3pt);
  \draw[arr, fmRed] (P) -- ++(40:0.7) node[above right, slbl]{$\sigma_{rr}$};
  \draw[arr, fmBlue] (P) -- ++(130:0.7) node[above left, slbl]{$\sigma_{\theta r}$};
  \node[slbl, below right] at (P) {$r$};
  \fill[fmBlue] (0,0) circle (2pt);
  \node[slbl, below left] at (0,0) {ось $z$ (линия)};
\end{tikzpicture}
\caption{}
\end{subfigure}
\caption{Компоненты напряжения: \textit{а}~— в декартовых координатах
(показаны компоненты на грани с нормалью $x$: нормальная $\sigma_{xx}$ и
касательные $\sigma_{yx}$, $\sigma_{zx}$); \textit{б}~— цилиндрические
координаты вокруг линии дислокации, направленной вдоль оси $z$}
\label{fig:stress-comp}
\end{figure}

\subsubsection{Поле напряжений винтовой дислокации}

Поле винтовой дислокации хорошо представить на модели упругого
цилиндрического кольца (рис.~\ref{fig:screw-cyl}). Сделаем в нём
радиальный разрез вдоль оси $z$ и сдвинем берега разреза друг
относительно друга на $b$ вдоль оси $z$. Внутреннюю область радиусом
$r_0$ (ядро) из рассмотрения исключаем.

\begin{figure}[ht]
\centering
\begin{tikzpicture}[scale=1.0, line join=round]
  \def\R{1.8}\def\r{0.45}\def\Hh{3.0}\def\e{0.35}\def\bb{0.4}
  % наружная поверхность
  \fill[fmBlue!10] (-\R,0) -- (-\R,\Hh) arc (180:360:\R cm and \e cm) -- (\R,0) arc (0:-180:\R cm and \e cm) -- cycle;
  \draw[atomLine] (-\R,0) -- (-\R,\Hh); \draw[atomLine] (\R,0) -- (\R,\Hh);
  \draw[atomLine] (\R,0) arc (0:-180:\R cm and \e cm);
  \draw[atomLine, dashed] (\R,0) arc (0:180:\R cm and \e cm);
  % верхний торец
  \filldraw[fill=fmBlue!5, draw=atomLine] (0,\Hh) ellipse (\R cm and \e cm);
  \filldraw[fill=white, draw=atomLine] (0,\Hh) ellipse (\r cm and 0.09cm);
  % разрез и сдвиг
  \fill[fmAmber!45] (\r,\Hh) -- (\R,\Hh) -- (\R,\Hh-\bb) -- (\r,\Hh-0.12) -- cycle;
  \draw[atomLine] (\r,\Hh) -- (\R,\Hh) (\r,\Hh-0.12) -- (\R,\Hh-\bb) -- (\R,0-\bb+0.0);
  \draw[bvec] (\R+0.3,\Hh) -- (\R+0.3,\Hh-\bb) node[right, slbl, fmRed]{$\vb b$};
  \draw[disl, dashed] (0,-0.6) -- (0,\Hh+0.7) node[above, slbl, fmBlue]{$z$, линия дислокации};
  \draw[dimline] (0,\Hh+0.25) -- node[above, slbl]{$r_0$} (\r,\Hh+0.25);
\end{tikzpicture}
\caption{Модель поля винтовой дислокации: упругое кольцо с радиальным
разрезом, берега которого сдвинуты на $b$ вдоль оси $z$. Ядро радиусом
$r_0$ исключено}
\label{fig:screw-cyl}
\end{figure}

\begin{shagi}
\item \emph{Симметрия.} Линия дислокации~— ось $z$, $\vb b=(0,0,b)$.
Поле обладает цилиндрической симметрией: все направления поперёк оси
равноправны. Поэтому $u_x=u_y=0$: отличные от нуля $u_x$ где-то
обязательно нарушили бы симметрию. Кроме того, ничего не меняется вдоль
оси: $\partial u_i/\partial z=0$.
\item \emph{Смещение.} Остаётся $u_z(r,\theta)$. На окружности
$r=\text{const}$ все точки равноправны, поэтому $u_z$ растёт
пропорционально углу; полный обход даёт $b$:
\begin{equation}\label{eq:screw-u}
u_z=\frac{b}{2\pi}\,\theta=\frac{b}{2\pi}\operatorname{arctg}\frac yx .
\end{equation}
\item \emph{Деформации.} Дифференцируя, получаем только две ненулевые
(сдвиговые) компоненты:
\[
\varepsilon_{xz}=-\frac{b}{4\pi}\,\frac{y}{x^2+y^2},\qquad
\varepsilon_{yz}=\frac{b}{4\pi}\,\frac{x}{x^2+y^2} .
\]
\item \emph{Напряжения} по закону Гука $\sigma_{ij}=2G\varepsilon_{ij}$:
\[
\sigma_{xz}=-\frac{Gb}{2\pi}\,\frac{y}{x^2+y^2},\qquad
\sigma_{yz}=\frac{Gb}{2\pi}\,\frac{x}{x^2+y^2} .
\]
\end{shagi}
В цилиндрических координатах всё ещё проще~— есть одна компонента:
\begin{keyeq}\label{eq:screw-sigma}
\sigma_{\theta z}=\sigma_{z\theta}=\frac{Gb}{2\pi r},\qquad\text{остальные}\ \sigma_{ij}=0 .
\end{keyeq}

\begin{glavnoe}
Поле винтовой дислокации~— \emph{чистый сдвиг}, убывающий как $1/r$.
Нормальных напряжений нет, поэтому нет и гидростатического давления
($p=-\sigma_{ll}/3=0$).
\end{glavnoe}

\subsubsection{Поле напряжений краевой дислокации}

Пусть линия краевой дислокации идёт вдоль $z$, а $\vb b=(b,0,0)$
(экстраплоскость~— полуплоскость $x=0$, $y>0$). Вдоль оси $z$ ничего не
меняется и смещения вдоль $z$ нет: деформация \emph{плоская}. Решение
теории упругости:
\begin{keyeq}\label{eq:edge-sigma}
\begin{aligned}
\sigma_{xx}&=-D\,\frac{y\,(3x^2+y^2)}{(x^2+y^2)^2}, &
\sigma_{yy}&=D\,\frac{y\,(x^2-y^2)}{(x^2+y^2)^2},\\
\sigma_{xy}&=D\,\frac{x\,(x^2-y^2)}{(x^2+y^2)^2}, &
\sigma_{zz}&=\nu(\sigma_{xx}+\sigma_{yy}),\qquad D=\frac{Gb}{2\pi(1-\nu)} .
\end{aligned}
\end{keyeq}
Остальные компоненты ($\sigma_{xz}$, $\sigma_{yz}$) равны нулю. В полярных
координатах:
\[
\sigma_{rr}=\sigma_{\theta\theta}=-D\,\frac{\sin\theta}{r},\qquad
\sigma_{r\theta}=D\,\frac{\cos\theta}{r},\qquad \sigma_{zz}=\nu(\sigma_{rr}+\sigma_{\theta\theta}).
\]

Почему $\sigma_{zz}\neq0$, хотя деформации вдоль $z$ нет? Растяжение или
сжатие по $x$ и $y$ вызвало бы пуассоново сокращение по $z$, но
соседние слои материала его не допускают~— отсюда напряжение $\sigma_{zz}$.

Особенно важно \term{гидростатическое давление} $p=-\frac13(\sigma_{xx}+
\sigma_{yy}+\sigma_{zz})$. Подставив~\eqref{eq:edge-sigma}, получим
\begin{equation}\label{eq:edge-p}
p=\frac{2(1+\nu)}{3}\,D\,\frac{y}{x^2+y^2}=\frac{2(1+\nu)}{3}\,D\,\frac{\sin\theta}{r}.
\end{equation}
Над плоскостью скольжения ($y>0$, там экстраплоскость) $p>0$~—
\emph{сжатие}, под ней~— \emph{растяжение} (рис.~\ref{fig:edge-field}).

\begin{figure}[ht]
\centering
\begin{tikzpicture}
\begin{axis}[width=0.56\textwidth, height=0.56\textwidth, view={0}{90},
  xmin=-3, xmax=3, ymin=-3, ymax=3, axis equal image,
  xlabel={$x/b$}, ylabel={$y/b$}, colorbar, colorbar style={width=3mm},
  colormap={rb}{rgb255(0cm)=(33,102,172); rgb255(1cm)=(146,197,222); rgb255(2cm)=(247,247,247); rgb255(3cm)=(244,165,130); rgb255(4cm)=(178,24,43)},
  point meta min=-1, point meta max=1]
  \addplot3[surf, shader=interp, domain=-3:3, domain y=-3:3, samples=61, samples y=61]
     {max(-1,min(1, 0.6*y/(x^2+y^2+0.05)))};
  \draw[black, line width=1.4pt] (axis cs:-0.3,0,1) -- (axis cs:0.3,0,1) (axis cs:0,0,1) -- (axis cs:0,0.45,1);
  \draw[black, dashed] (axis cs:-3,0,1) -- (axis cs:3,0,1);
  \node[font=\small, fill=white, fill opacity=0.7, text opacity=1] at (axis cs:0,2.3,1) {сжатие};
  \node[font=\small, fill=white, fill opacity=0.7, text opacity=1] at (axis cs:0,-2.3,1) {растяжение};
\end{axis}
\end{tikzpicture}
\caption{Гидростатическое давление вокруг краевой дислокации ($\perp$ в
центре). Над плоскостью скольжения (пунктир)~— сжатие (красный), под
ней~— растяжение (синий). Линии равного давления~— окружности, касающиеся
плоскости скольжения в точке выхода дислокации}
\label{fig:edge-field}
\end{figure}

\paragraph{Сравнение полей}
\begin{table}[ht]
\centering
\caption{Сравнение полей напряжений винтовой и краевой дислокаций}
\label{tab:field-compare}
\small
\begin{tabular}{@{}lcc@{}}
\toprule
 & Винтовая & Краевая \\
\midrule
Ненулевые компоненты & $\sigma_{xz}$, $\sigma_{yz}$ & $\sigma_{xx}$, $\sigma_{yy}$, $\sigma_{zz}$, $\sigma_{xy}$ \\
Деформация & чистый сдвиг & плоская деформация \\
Гидростатическое давление & нет & есть ($\propto\sin\theta/r$) \\
Убывание с расстоянием & $1/r$ & $1/r$ \\
\bottomrule
\end{tabular}
\end{table}
Отсутствие гидростатического давления у винтовой дислокации приводит к
коренным различиям во взаимодействии краевых и винтовых дислокаций с
точечными дефектами (см. ниже).

\paragraph{Влияние поверхности}
В бесконечной среде поле $\propto1/r$ не обращается в нуль ни на каком
расстоянии. В реальном кристалле на свободной поверхности нормальные
напряжения должны исчезать; чтобы это учесть, к полю бесконечной среды
добавляют поле поверхностных сил обратного знака. Поправка существенна
лишь там, где до поверхности ближе, чем до оси дислокации; обычно
пользуются формулами для бесконечной среды. Однако вблизи поверхности
поправка важна: энергия поля убывает при приближении дислокации к
поверхности (часть поля «заменяется пустотой»), поэтому
\vazhno{дислокация притягивается к свободной поверхности}. Эту силу
считают, как в электростатике: как притяжение к фиктивной
«дислокации-изображению» противоположного знака, расположенной
зеркально относительно поверхности.

\paragraph{Где кончается теория упругости}
Все формулы дают $\sigma\to\infty$ при $r\to0$~— это, конечно,
неверно. Деформации $\varepsilon\sim b/(4\pi r)$ малы ($\lesssim0{,}05$)
лишь при $r\gtrsim(1{,}5\text{--}3)\,b$; ближе работает не упругость, а
дискретная атомная модель ядра. Её «сшивают» с упругим полем на границе
ядра, чтобы найти энергию ядра и напряжение Пайерлса.

\subsubsection{Энергия дислокации}

Вокруг дислокации есть поле напряжений, значит, с ней связана энергия
искажения. Она равна работе, затраченной на создание дислокации, и
складывается из упругой энергии и энергии ядра:
$E=E_{\text{упр}}+E_{\text{ядр}}$.

\paragraph{Винтовая дислокация}
Создадим дислокацию единичной длины так, как на рис.~\ref{fig:screw-cyl}:
будем постепенно сдвигать берега разреза до смещения $b$. На площадке
разреза шириной $\d r$ (на расстоянии $r$ от оси) действует напряжение,
которое по мере сдвига линейно растёт от нуля до $\sigma_{\theta z}$;
среднее значение $\frac12\sigma_{\theta z}$. Работа на этой площадке
$\frac12\sigma_{\theta z}\,b\,\d r$. Суммируя от ядра $r_0$ до внешнего
радиуса $R$:
\[
E_{\text{упр}}=\int_{r_0}^{R}\frac12\,\frac{Gb}{2\pi r}\,b\,\d r=\frac{Gb^2}{4\pi}\ln\frac{R}{r_0} .
\]

\paragraph{Краевая и смешанная дислокации}
Аналогичный расчёт для краевой дислокации (разрез по плоскости
скольжения, где действует $\sigma_{r\theta}=D/r$) даёт ответ, больший в
$1/(1-\nu)\approx1{,}5$ раза. Смешанную дислокацию, у которой $\vb b$
составляет угол $\theta$ с линией, можно представить как сумму краевой
(вектор $b\sin\theta$) и винтовой ($b\cos\theta$) дислокаций:
\begin{keyeq}\label{eq:E-disl}
E_{\text{упр}}=\frac{Gb^2}{4\pi(1-\nu)}\bigl(1-\nu\cos^2\theta\bigr)\ln\frac{R}{r_0}\approx\alpha\,Gb^2,\qquad\alpha\approx0{,}5\text{--}1 .
\end{keyeq}
При $\theta=0$ это энергия винтовой, при $\theta=90^\circ$~— краевой
дислокации.

\paragraph{Чему равен радиус обрезания $R$}
В бесконечном кристалле энергия логарифмически расходится. В реальном
кристалле поля соседних дислокаций местами усиливают, местами гасят друг
друга, и поле данной дислокации в среднем компенсируется на расстоянии
порядка расстояния между дислокациями: $R\approx\rho^{-1/2}$, где
$\rho$~— плотность дислокаций. Логарифм меняется слабо: при
$\rho=10^8$~см$^{-2}$ и $r_0=10^{-7}$~см $\ln(R/r_0)=\ln10^3\approx7$.

\begin{primer}[Энергия дислокации]
Краевая дислокация, $G=40$~ГПа ($4\cdot10^{11}$~дин/см$^2$), $b=2{,}5\cdot
10^{-8}$~см, $\nu=0{,}3$, $R=1$~см, $r_0=10^{-7}$~см:
\[
E_{\text{упр}}=\frac{4\cdot10^{11}\cdot(2{,}5\cdot10^{-8})^2}{4\pi\cdot0{,}7}\ln10^7
\approx5\cdot10^{-4}~\text{эрг/см}.
\]
На отрезок длиной в одно межатомное расстояние ($b$) приходится
$5\cdot10^{-4}\cdot2{,}5\cdot10^{-8}\approx1{,}2\cdot10^{-11}$~эрг
$\approx7$--$8\eV$\ispr{В оригинале (с.~238) размерность энергии
дислокации указана как «эрг$\cdot$см$^{-2}$»; энергия на единицу длины
линии измеряется в эрг/см (Дж/м). Пропущен и порядок модуля сдвига
($G=4\cdot10^{11}$~дин/см$^2$).}. Энергия ядра оценивается в
$0{,}5$--$1{,}5\eV$ на атомную плоскость~— лишь небольшая доля упругой.
\end{primer}

\begin{glavnoe}
\begin{enumerate}[label=\arabic*)]
\item Энергия дислокации \vazhno{пропорциональна $b^2$}: выгодны
дислокации с наименьшим вектором Бюргерса (это основа критерия Франка,
п.~\ref{ssec:disl-in-structures}).
\item Энергия \vazhno{пропорциональна длине}: дислокация стремится
сократиться, как натянутая струна. Её \term{линейное натяжение}
$T\approx\alpha Gb^2\approx\frac12Gb^2$.
\item Энергии краевой, винтовой и смешанной дислокаций одного порядка.
\end{enumerate}
\end{glavnoe}

\paragraph{Дислокации~— неравновесные дефекты}
Энергия даже короткого отрезка дислокации длиной $b$ велика:
$E\approx\frac12Gb^3\approx2{,}5\eV$ (при $G=30$~ГПа, $b=3\cdot10^{-8}$~см).
Оценим равновесную концентрацию самой маленькой дислокационной петли
радиусом $3b$ при комнатной температуре. Её энергия
$\frac12Gb^3\cdot\frac{2\pi\cdot3b}{b}\approx2{,}5\cdot19\approx47\eV$, и
даже если число возможных мест $N_0=10^{22}$~см$^{-3}$ (не больше числа
атомов),
\[
N=N_0\Ex{-\frac{E}{kT}}=10^{22}\exp\left(-\frac{47}{0{,}025}\right)\sim10^{-800}~\text{см}^{-3}.
\]
Одна петля на объём, во много раз превосходящий видимую Вселенную!

\begin{glavnoe}
Дислокации никогда не бывают термодинамически равновесными дефектами
(в отличие от вакансий). Равновесию отвечает нулевая плотность
дислокаций. Дислокации~— следствие предыстории кристалла: роста,
деформации, термообработки. Но из-за высоких энергетических барьеров
полное равновесие недостижимо, и никакой отжиг не уменьшает плотность
дислокаций до нуля.
\end{glavnoe}

\subsubsection{Силы, действующие на дислокацию}

Движение дислокации вызывает пластическую деформацию, т.\,е. совершает
работу. Поэтому удобно говорить о силе, действующей на дислокацию, как
на самостоятельный объект.

\begin{vnimanie}
Дислокация~— не частица, а особая конфигурация расположения атомов.
«Сила, действующая на дислокацию»,~— это не сила, приложенная к каким-то
атомам, а \emph{конфигурационная} (виртуальная) сила: изменение
свободной энергии системы при перемещении дислокации на единицу длины.
Вспомните винтовую дислокацию: атомы её ядра смещаются вдоль $\vb b$ под
действием приложенных к ним сил, а сама дислокация движется
перпендикулярно им.
\end{vnimanie}

\paragraph{Сила скольжения}
Пусть в кристалле размером $L_1\times L_2$ в плоскости скольжения
касательное напряжение $\tau$ действует в направлении $\vb b$. Когда
дислокация длиной $L_1$ пройдёт всю плоскость (путь $L_2$), верхняя часть
кристалла сдвинется на $b$, и внешняя сила $\tau L_1L_2$ совершит работу
$\tau L_1L_2b$. Ту же работу совершает сила $F$, действующая на единицу
длины дислокации, на пути $L_2$: $FL_1L_2$. Отсюда
\begin{keyeq}\label{eq:F-taub}
F=\tau b .
\end{keyeq}
Сила перпендикулярна линии дислокации и лежит в плоскости скольжения
(даже для винтовой дислокации, у которой $\vb b$ параллелен линии!).

\paragraph{Формула Пича--Кёлера}
В общем случае, когда тензор напряжений $\tn\sigma$ произволен, сила на
единицу длины дислокации
\begin{keyeq}\label{eq:peach-koehler}
\vb f=(\tn\sigma\cdot\vb b)\times\vb t .
\end{keyeq}
Здесь $\tn\sigma\cdot\vb b$~— вектор с компонентами $\sigma_{ij}b_j$
(сила, действующая на площадку с нормалью $\vb b$, умноженная на $b$).
Знак определяется соглашением о направлениях $\vb t$ и $\vb b$ (правило
FS/RH; при другом соглашении $\vb f=\vb t\times(\tn\sigma\cdot\vb b)$).

\begin{podrobnee}[Вывод формулы Пича--Кёлера]
Пусть отрезок дислокации единичной длины смещается на $\d\vb s$ и заметает
площадку $\d\vb A=\vb t\times\d\vb s$ (вектор, нормальный к площадке).
Материал по разные стороны площадки смещается друг относительно друга на
$\vb b$. Сила, действующая на площадку $\d\vb A$, равна $\tn\sigma\cdot
\d\vb A$, и работа напряжений
\[
\d W=\vb b\cdot(\tn\sigma\cdot\d\vb A)=(\tn\sigma\cdot\vb b)\cdot(\vb t\times\d\vb s)
=\bigl[(\tn\sigma\cdot\vb b)\times\vb t\bigr]\cdot\d\vb s ,
\]
где использована симметрия тензора напряжений и свойство смешанного
произведения. Сравнивая с $\d W=\vb f\cdot\d\vb s$, получаем
$\vb f=(\tn\sigma\cdot\vb b)\times\vb t$.
\end{podrobnee}

\begin{primer}[Силы на краевую дислокацию]
Пусть линия дислокации идёт вдоль оси $x_2$, $\vb b=(b,0,0)$, плоскость
скольжения~— $x_1x_2$ (нормаль $x_3$). Тогда
$\tn\sigma\cdot\vb b=b(\sigma_{11},\sigma_{21},\sigma_{31})$ и
\[
\vb f=b(\sigma_{11},\sigma_{21},\sigma_{31})\times(0,1,0)=b\,(-\sigma_{31},\,0,\,\sigma_{11}).
\]
\begin{itemize}
\item Компонента вдоль $x_1$ (в плоскости скольжения)~— \emph{сила
скольжения} $b\sigma_{31}$: её создаёт касательное напряжение в
плоскости скольжения вдоль $\vb b$, как и в~\eqref{eq:F-taub}.
\item Компонента вдоль $x_3$ (перпендикулярно плоскости скольжения)~—
\emph{сила переползания} $b\sigma_{11}$: её создаёт \emph{нормальное}
напряжение вдоль $\vb b$. Сжатие стремится «выдавить» экстраплоскость
из кристалла, растяжение~— «втянуть» её.
\item Напряжение $\sigma_{21}$ даёт силу вдоль самой линии и работы не
совершает.
\end{itemize}
\end{primer}

\paragraph{«Химическая» сила}
Сила переползания уравновешивается, когда концентрация вакансий у
дислокации равна равновесной. Если же она отличается (пересыщение
$C/C_0$), на дислокацию действует осмотическая сила (п.~\ref{ssec:disl-motion})
$f\approx(kT/b^2)\ln(C/C_0)$. При комнатной температуре ($kT=1/40\eV$,
$Gb^3\approx5\eV$) соответствующее «химическое напряжение»
\[
\frac fb\approx\frac{G}{200}\ln\frac{C}{C_0} .
\]
Немногие кристаллы выдерживают механические напряжения больше $G/200$,
поэтому даже небольшие пересыщения вакансиями, например после закалки,
вызывают заметное переползание~— скажем, рост мелких призматических
петель (это наблюдали в алюминии).

\subsubsection{Взаимодействие дислокаций}

\paragraph{Качественная картина}
Рассмотрим две параллельные краевые дислокации с одинаковым $b$ в одной
плоскости скольжения. Пока они далеко, их энергия равна сумме энергий:
$2\cdot\alpha Gb^2$.
\begin{itemize}
\item Если дислокации \emph{одного знака} слить, получится дислокация с
вектором $2\vb b$ и энергией $\alpha G(2b)^2=4\alpha Gb^2$~— вдвое больше.
Значит, дислокации одного знака \vazhno{отталкиваются}.
\item Если \emph{разного знака}, суммарный вектор Бюргерса равен нулю, и
при слиянии энергия поля исчезает. Значит, они \vazhno{притягиваются} и
аннигилируют (если лежат в разных, близких плоскостях~— с образованием
ряда вакансий или межузельных атомов, рис.~\ref{fig:annih}).
\end{itemize}

\paragraph{Две параллельные краевые дислокации}
Поместим дислокацию~2 в поле дислокации~1, находящейся в начале координат.
По формуле Пича--Кёлера компоненты силы на единицу длины дислокации~2:
\begin{equation}\label{eq:edge-edge}
f_x=\sigma_{xy}\,b=\frac{Gb^2}{2\pi(1-\nu)}\cdot\frac{x\,(x^2-y^2)}{(x^2+y^2)^2},\qquad
f_y=-\sigma_{xx}\,b=\frac{Gb^2}{2\pi(1-\nu)}\cdot\frac{y\,(3x^2+y^2)}{(x^2+y^2)^2}.
\end{equation}
Краевая дислокация легко движется только скольжением, поэтому главная~—
компонента $f_x$ (рис.~\ref{fig:edge-edge}).

\begin{figure}[ht]
\centering
\begin{subfigure}[b]{0.56\textwidth}\centering
\begin{tikzpicture}
\begin{axis}[width=\linewidth, height=5.2cm, xmin=-3, xmax=3, ymin=-0.3, ymax=0.3,
  axis lines=middle, xlabel={$x/y$}, ylabel={$f_x$},
  xtick={-2,-1,1,2}, ytick=\empty,
  every axis x label/.style={at={(ticklabel* cs:1)}, anchor=west},
  every axis y label/.style={at={(ticklabel* cs:1)}, anchor=south}]
  \addplot[fmBlue, very thick, domain=-3:3, samples=200] {x*(x^2-1)/(x^2+1)^2};
  \addplot[only marks, mark=*, mark size=2.2pt, fmGreen] coordinates {(0,0)};
  \addplot[only marks, mark=o, mark size=2.2pt, fmRed, thick] coordinates {(-1,0) (1,0)};
  \node[font=\footnotesize, anchor=west, align=left] at (axis cs:-2.95,0.22) {\textcolor{fmGreen}{$\bullet$}~устойчиво\\\textcolor{fmRed}{$\circ$}~неустойчиво};
\end{axis}
\end{tikzpicture}
\caption{}
\end{subfigure}\hfill
\begin{subfigure}[b]{0.4\textwidth}\centering
\begin{tikzpicture}[scale=0.85]
  \draw[fmGray, dashed] (-0.6,0) -- (0.6,0) (-0.6,1) -- (0.6,1) (-0.6,2) -- (0.6,2);
  \foreach \y in {0,1,2}{\dislT[fmBlue]{(0,\y)}{0}}
  \node[slbl, align=center] at (0,-0.7) {одного знака:\\стенка};
  \draw[fmGray, dashed] (2.0,0.5) -- (3.4,0.5) (2.4,1.5) -- (3.8,1.5);
  \dislT[fmBlue]{(2.7,0.5)}{0}
  \dislT[fmRed]{(3.7,1.5)}{180}
  \draw[fmGray, dotted] (2.7,0.5) -- (3.7,1.5);
  \node[slbl, align=center] at (3.2,-0.7) {разного знака:\\диполь под $45^\circ$};
\end{tikzpicture}
\caption{}
\end{subfigure}
\caption{Взаимодействие параллельных краевых дислокаций: \textit{а}~—
сила $f_x$ на дислокацию~2 одного знака с дислокацией~1, как функция
$x$ при фиксированном расстоянии $y$ между плоскостями скольжения (для
дислокаций разного знака знак силы меняется); \textit{б}~— устойчивые
конфигурации}
\label{fig:edge-edge}
\end{figure}

\begin{itemize}
\item $f_x=0$ при $x=0$ и при $x=\pm y$.
\item Для дислокаций \emph{одного знака} положение $x=0$ (одна над
другой) \emph{устойчиво}, а $x=\pm y$~— неустойчиво. Поэтому одноимённые
краевые дислокации выстраиваются друг над другом, образуя
\term{дислокационную стенку} (п.~\ref{ssec:disl-structures}).
\item Для дислокаций \emph{разного знака} всё наоборот: устойчивы
положения $x=\pm y$ (линия, соединяющая дислокации, наклонена под
$45^\circ$)~— это \term{дислокационный диполь}.
\item Если дислокации лежат в одной плоскости скольжения ($y=0$):
$f_x=\dfrac{Gb_1b_2}{2\pi(1-\nu)\,x}$~— отталкивание одноимённых и
притяжение разноимённых, убывающее как $1/x$.
\end{itemize}

\paragraph{Две параллельные винтовые дислокации}
Поле винтовой дислокации осесимметрично, и сила направлена по линии,
соединяющей дислокации:
\begin{equation}\label{eq:screw-screw}
f_r=\sigma_{\theta z}b=\frac{Gb^2}{2\pi r}.
\end{equation}
Одноимённые винтовые дислокации отталкиваются, разноимённые~—
притягиваются. В общем случае сила, действующая на дислокацию~2 в поле
дислокации~1, вычисляется по той же формуле Пича--Кёлера:
$\vb f_2=(\tn\sigma_1\cdot\vb b_2)\times\vb t_2$.

\subsubsection{Взаимодействие дислокаций с точечными дефектами}

Точечный дефект искажает решётку вокруг себя и поэтому взаимодействует с
полем напряжений дислокации. Если дефект подходит к дислокации (или
она к нему), энергия искажений системы может понизиться. Чтобы затем
оторвать дислокацию от дефекта, нужна дополнительная работа~— кристалл
\emph{упрочняется}.

\paragraph{Размерное взаимодействие}
Пусть атом примеси радиусом $R(1+\varepsilon)$ замещает атом матрицы
радиусом $R$. Он раздвигает окружающий материал~— объём меняется на
$\Delta V=4\pi\varepsilon R^3$ (с точностью до малых высших степеней
$\varepsilon$). Если в этом месте действует гидростатическое давление
$p$, энергия взаимодействия равна работе против него:
$E_{\text{вз}}=p\,\Delta V$. Подставляя давление краевой
дислокации~\eqref{eq:edge-p}:
\begin{keyeq}\label{eq:E-int}
E_{\text{вз}}=\frac{4(1+\nu)}{3(1-\nu)}\,Gb\,\varepsilon R^3\,\frac{y}{x^2+y^2}
=A\,\frac{\sin\theta}{r},\qquad A=\frac{4(1+\nu)}{3(1-\nu)}\,Gb\,\varepsilon R^3 .
\end{keyeq}

\begin{figure}[ht]
\centering
\begin{tikzpicture}[scale=0.85]
  \foreach \d/\o in {0.6/70, 1.1/50, 1.8/32, 2.7/20}{
    \draw[fmRed!\o!white, line width=0.8pt] (0,\d/2) circle (\d/2);
    \draw[fmBlue!\o!white, line width=0.8pt] (0,-\d/2) circle (\d/2);}
  \draw[fmGray, dashed] (-3.2,0) -- (3.2,0);
  \dislT[black]{(0,0)}{0}
  \node[slbl, fmRed, align=left, right] at (1.6,1.9) {отталкивание\\($E_{\text{вз}}>0$)};
  \node[slbl, fmBlue, align=left, right] at (1.6,-1.9) {притяжение\\($E_{\text{вз}}<0$)};
  \node[atomB=3.4mm] at (0,-0.38) {};
\end{tikzpicture}
\caption{Линии равной энергии взаимодействия краевой дислокации с крупным
атомом примеси ($\varepsilon>0$)~— окружности, касающиеся плоскости
скольжения. Чем меньше окружность, тем больше $|E_{\text{вз}}|$. Атом
выгоднее всего располагать прямо под краем экстраплоскости}
\label{fig:cottrell-iso}
\end{figure}

\begin{itemize}
\item Крупный атом ($\varepsilon>0$) отталкивается от сжатой области и
притягивается к растянутой; выгоднее всего ему сидеть непосредственно под
краем экстраплоскости ($\theta=3\pi/2$, малое $r$). Мелкий атом
($\varepsilon<0$)~— наоборот, над краем.
\item Разность энергий атома у дислокации и вдали от неё~—
\term{энергия связи} дефекта с дислокацией. Для краевой дислокации эта
простая модель даёт реалистичные значения $\sim0{,}1$--$1\eV$.
\item Линии $E_{\text{вз}}=\text{const}$, т.\,е. $\sin\theta/r=\text{const}$,~—
окружности, касающиеся плоскости скольжения в точке выхода дислокации
(рис.~\ref{fig:cottrell-iso}).
\end{itemize}

\paragraph{А винтовая дислокация?}
У винтовой дислокации нет гидростатического давления, и со «сферическим»
дефектом она по этому механизму не взаимодействует ($E_{\text{вз}}=0$).
Однако если дефект искажает решётку \emph{несимметрично}, он чувствует и
касательные напряжения. Так, атомы C и N в октаэдрических междоузлиях
ОЦК-железа (в центрах граней и серединах рёбер куба) создают
\emph{тетрагональное} искажение. Такие атомы взаимодействуют и с краевыми,
и с винтовыми дислокациями. Для винтовой дислокации вдоль $\nfam{111}$:
\[
U_{\text{вз}}=\frac{2}{3\pi}\,GbR^3(\varepsilon_1-\varepsilon_2)\,\frac{\cos\theta}{r},
\]
где $\varepsilon_1$~— деформация вдоль тетрагональной оси, $\varepsilon_2$~—
поперёк. Для углерода в железе $\varepsilon_1=0{,}38$,
$\varepsilon_2=-0{,}026$, и энергия связи достигает
$\sim0{,}75\eV$\ispr{В оригинале (с.~249) указано $\varepsilon_2=0{,}26$.
В классической работе Кохардта, Шёка и Видерзиха, откуда взяты эти
значения, $\varepsilon_2=-0{,}026$ (сжатие поперёк тетрагональной оси).}.
Смешанная дислокация притягивает любые атомы, поскольку у неё есть
краевая компонента: чем ближе угол между линией и $\vb b$ к прямому, тем
сильнее притяжение.

\paragraph{Атмосферы Коттрелла}
\begin{opredelenie}
\term{Атмосфера}~— неоднородное распределение примесных атомов вокруг
дислокации, вызванное их взаимодействием. Различают атмосферы
\term{Коттрелла} (упругое размерное взаимодействие с краевой
дислокацией), \term{Снука} (упорядочение атомов внедрения в поле
напряжений) и \term{Сузуки} (химическое взаимодействие с дефектом
упаковки).
\end{opredelenie}

В равновесии концентрация примеси в точке, где энергия связи равна $E$
(здесь $E>0$ означает притяжение), подчиняется распределению Больцмана:
\begin{equation}\label{eq:cottrell}
C=C_0\Ex{\frac{E}{kT}},
\end{equation}
где $C_0$~— средняя концентрация примеси в металле. Поэтому линии равной
энергии (рис.~\ref{fig:cottrell-iso}) одновременно являются линиями равной
концентрации. С повышением температуры атмосфера «рассасывается», с
понижением~— густеет, и вблизи ядра могут даже образоваться мелкие
выделения второй фазы.

\paragraph{Насыщенная атмосфера}
В местах с наибольшей энергией связи $E_{\max}$ (под краем
экстраплоскости) концентрация $C_{\max}=C_0\exp(E_{\max}/kT)$ при
охлаждении растёт и наконец достигает единицы: все такие места вдоль
дислокации заняты. Атмосферу называют \term{насыщенной}
(конденсированной). Из условия $C_{\max}=1$ находим температуру, ниже
которой атмосфера насыщена:
\begin{keyeq}\label{eq:T-sat}
T_\text{н}=\frac{E_{\max}}{k\ln(1/C_0)} .
\end{keyeq}

\begin{primer}[Насыщение атмосфер]
\begin{itemize}
\item Углерод и азот в $\alpha$-железе (атомы внедрения), $E_{\max}\approx0{,}5\eV$,
$C_0=10^{-4}$: $T_\text{н}=0{,}5/(8{,}6\cdot10^{-5}\cdot9{,}2)\approx630$~К.
При комнатной температуре атмосферы насыщены.
\item Сплав меди с 1~ат.\,\% цинка (атомы замещения, ГЦК), $E_{\max}\approx0{,}1\eV$:
$T_\text{н}=0{,}1/(8{,}6\cdot10^{-5}\cdot4{,}6)\approx250$--$300$~К. Малую
энергию связи компенсирует большая концентрация.
\end{itemize}
Сколько примеси нужно, чтобы «одеть» все дислокации одноатомными
цепочками? Доля таких атомов $C=\rho a^2$: при $\rho=10^8$~см$^{-2}$ (отожжённый
металл) $C\sim10^{-7}$, при $\rho=10^{11}$~см$^{-2}$ (наклёп) $C\sim10^{-4}$.
Даже в металле технической чистоты примесей хватает.
\end{primer}

Энергия связи атомов внедрения с дислокацией значительно выше, чем атомов
замещения. Поэтому при одинаковой концентрации $C_0$ в растворах
внедрения $T_\text{н}$ выше.

\paragraph{Отрыв дислокации от атмосферы}
Любая атмосфера тормозит движение дислокаций. Найдём силу, нужную, чтобы
оторвать краевую дислокацию от ряда примесных атомов, сидящих на
расстоянии $r_0$ под плоскостью скольжения (рис.~\ref{fig:unpin}).

\begin{shagi}
\item Когда дислокация смещена по плоскости скольжения на $x$,
расстояние до ряда $r=\sqrt{x^2+r_0^2}$, $\sin\theta=r_0/r$, и энергия
связи на одну атомную плоскость $E=Br_0/(x^2+r_0^2)$ (постоянная $B$~—
это $A$ из~\eqref{eq:E-int} в расчёте на атомную плоскость).
\item Сила, удерживающая дислокацию: $F_x=\dfrac{\partial E}{\partial x}$,
по модулю $\dfrac{2Br_0x}{(x^2+r_0^2)^2}$.
\item Максимум силы (приравниваем производную нулю) достигается при
$x=r_0/\sqrt3$:
\end{shagi}
\begin{equation}\label{eq:Fmax-unpin}
F_{\max}=\frac{3\sqrt3}{8}\,\frac{B}{r_0^2}.
\end{equation}
Это и есть сила (на атомную плоскость), которую нужно приложить, чтобы
оторвать дислокацию от насыщенной атмосферы. С этим связан эффект резкого
предела текучести в малоуглеродистых сталях.

\begin{figure}[ht]
\centering
\begin{tikzpicture}
\begin{axis}[width=0.6\textwidth, height=4.8cm, xmin=-3, xmax=3, ymin=-1.1, ymax=0.75,
  axis lines=middle, xlabel={$x/r_0$}, ytick=\empty,
  xtick={-1.732,-0.577,0.577,1.732}, xticklabels={,$-\tfrac{1}{\sqrt3}$,$\tfrac{1}{\sqrt3}$,},
  every axis x label/.style={at={(ticklabel* cs:1)}, anchor=west},
  legend style={at={(0.02,0.98)}, anchor=north west}, legend cell align=left]
  \addplot[fmBlue, very thick, domain=-3:3, samples=150] {-1/(x^2+1)};
  \addlegendentry{энергия связи $-E(x)$}
  \addplot[fmRed, very thick, domain=-3:3, samples=150] {2*x/(x^2+1)^2};
  \addlegendentry{удерживающая сила}
  \draw[dashed, fmGray] (axis cs:0.577,0) -- (axis cs:0.577,0.6495);
\end{axis}
\end{tikzpicture}
\caption{Энергия и сила взаимодействия дислокации с рядом примесных атомов
при её смещении по плоскости скольжения на $x$ (в единицах $B/r_0$ и
$B/r_0^2$). Наибольшая удерживающая сила~— при $x=r_0/\sqrt3$}
\label{fig:unpin}
\end{figure}

\paragraph{Атмосферы Снука}
В ОЦК-железе атомы углерода и азота занимают октаэдрические междоузлия
трёх типов~— у каждого типа своя ось тетрагонального искажения ($x$,
$y$ или $z$). Без напряжений все три типа заселены одинаково. Если
кристалл растянуть вдоль одной из осей, междоузлия этой ориентации
немного «раздвинутся», и атомы внедрения начнут преимущественно занимать
их (\term{эффект Снука}). Такое же упорядочение происходит в поле
напряжений краевой и винтовой дислокаций; область упорядоченных атомов
внедрения вокруг дислокации~— \term{атмосфера Снука}. В отличие от
медленного образования атмосферы Коттрелла (нужна диффузия на большие
расстояния), упорядочение Снука возникает быстро~— за счёт перескоков
атомов в \emph{соседние} междоузлия.

\paragraph{Атмосферы Сузуки}
В ГЦК-решётке расщеплённая дислокация (п.~\ref{ssec:disl-in-structures})
содержит полосу дефекта упаковки~— тонкую прослойку с чередованием слоёв
ГПУ-структуры. Растворимость примеси в ГЦК- и ГПУ-прослойках, вообще
говоря, разная, и при достаточно высокой температуре примесь диффузионно
перераспределяется между ними, как между двумя фазами. Изменённую
концентрацию примеси в дефекте упаковки называют \term{атмосферой
Сузуки}, а само явление~— \emph{химическим взаимодействием} расщеплённой
дислокации с растворёнными атомами.

\paragraph{Вакансии и собственные межузельные атомы}
Межузельный атом (он «распирает» решётку) притягивается к растянутой
области краевой дислокации, вакансия~— к сжатой; упругое взаимодействие
вакансий слабее. Принципиальное отличие от примеси: вакансии и
межузельные атомы, подойдя к дислокации, могут \emph{исчезнуть} на
порогах (дислокация при этом переползает), а примесные атомы сохраняют
индивидуальность и образуют атмосферы. Лишь если дефекты подходят к
дислокации быстрее, чем исчезают на порогах, они тоже могут образовать
подобие атмосферы.

\paragraph{Электрическое взаимодействие}
Область растяжения у края экстраплоскости заряжена слегка отрицательно
(туда «перетекают» электроны), область сжатия~— положительно. Краевая и
смешанная дислокации~— слабые линейные электрические диполи и
взаимодействуют с заряженными примесными атомами. В металлах это
взаимодействие значительно слабее упругого (табл.~\ref{tab:binding}).
Вклад неупругого взаимодействия в ядре количественно не оценён: ошибка
упругих расчётов~— около 10\,\% на расстоянии двух межатомных расстояний
от ядра и около 50\,\% на расстоянии одного.

\begin{table}[ht]
\centering
\caption{Энергия связи точечных дефектов с краевой дислокацией (дефект на
расстоянии одного межатомного расстояния от ядра), эВ}
\label{tab:binding}
\small
\begin{tabular}{@{}lcc@{}}
\toprule
Точечный дефект & упругая & электрическая \\
\midrule
Межузельный атом & 0,2--0,5 & 0,02 \\
Атом замещения & 0,05--0,1 & 0,02 \\
Вакансия & 0,02 & 0,02 \\
\bottomrule
\end{tabular}
\end{table}

\begin{itogi}
\begin{itemize}
\item Поле дислокации убывает как $1/r$. Винтовая: чистый сдвиг
$\sigma_{\theta z}=Gb/2\pi r$, давления нет. Краевая: $D=Gb/2\pi(1-\nu)$,
над плоскостью скольжения сжатие, под ней растяжение.
\item Энергия $E\approx\frac{Gb^2}{4\pi K}\ln\frac{R}{r_0}\approx\alpha
Gb^2$ на единицу длины; линейное натяжение $T\approx\frac12Gb^2$.
Дислокации термодинамически неравновесны.
\item Сила Пича--Кёлера $\vb f=(\tn\sigma\cdot\vb b)\times\vb t$; сила
скольжения $\tau b$.
\item Одноимённые дислокации отталкиваются, разноимённые притягиваются;
одноимённые краевые устойчивы друг над другом (стенка).
\item Примеси притягиваются к дислокациям ($E=A\sin\theta/r$) и образуют
атмосферы Коттрелла (насыщенные ниже $T_\text{н}=E_{\max}/k\ln(1/C_0)$),
Снука, Сузуки; атмосферы закрепляют дислокации.
\end{itemize}
\end{itogi}

\begin{proverka}
\item Почему винтовая дислокация не взаимодействует с вакансией по
размерному механизму?
\item Во сколько раз энергия краевой дислокации больше энергии
винтовой при $\nu=1/3$?
\item Почему дислокации не могут быть равновесными дефектами?
\item Какие компоненты напряжения вызывают скольжение и переползание
краевой дислокации?
\item Почему одноимённые краевые дислокации выстраиваются в вертикальную
стенку?
\item Оцените $T_\text{н}$ для атмосфер из атомов замещения с
$E_{\max}=0{,}2\eV$ при $C_0=10^{-3}$.
\end{proverka}
